\chapter{Detection, Monitoring, and Security Operations in the Age of AI}

\begin{learningobjectives}
By the end of this chapter, the reader should be able to:
\begin{itemize}
    \item explain how a financial institution converts raw security telemetry into alerts, cases, incidents, and management decisions;
    \item understand why detection quality is an operational-risk variable rather than merely a technical metric;
    \item describe how frontier AI models are changing vulnerability discovery, code analysis, reconnaissance, triage, and long-horizon cyber operations;
    \item recognize, at a safe and non-operational level, the main AI-enabled tactics used to scale or automate cyber activity;
    \item explain how autonomous and multi-agent architectures can increase persistence, parallelism, and adaptability;
    \item design a defensive architecture that combines deterministic controls, AI-assisted analysis, human judgment, and governed autonomous response;
    \item connect security-operations performance to trading continuity, payments, client protection, market infrastructure, liquidity, and financial loss.
\end{itemize}
\end{learningobjectives}

\section{Security Operations as a Financial Control Function}

A security operations center, usually abbreviated SOC, is often presented as a technical room filled with screens. That description misses its economic purpose.

The SOC exists to answer a business question:

\begin{quote}
Is something occurring in our digital environment that could impair the institution's ability to operate safely, honor obligations, protect assets, or preserve trust?
\end{quote}

For a financial institution, this question is inseparable from operational risk.

A bank may have millions of authentication events each day. A broker-dealer may process large volumes of trading messages. An asset manager may receive prices, portfolio files, client instructions, and vendor data. An exchange may handle continuous participant connections and order flow. The defensive platform must distinguish ordinary activity from signals that justify investigation.

The monitoring pipeline can be represented as:

\begin{equation}
\text{Collect}
\rightarrow
\text{Normalize}
\rightarrow
\text{Enrich}
\rightarrow
\text{Detect}
\rightarrow
\text{Triage}
\rightarrow
\text{Escalate}
\rightarrow
\text{Respond}.
\end{equation}

Each stage matters.

Collection determines whether the institution has evidence.

Normalization makes evidence comparable.

Enrichment adds business context.

Detection identifies suspicious patterns.

Triage prioritizes analyst attention.

Escalation connects technical evidence to management.

Response limits the damage.

A weakness anywhere in this chain can increase the final loss.

\section{Why Detection Speed Matters in Finance}

Time has economic value in financial markets.

If an unauthorized user accesses a payment system for thirty seconds and is blocked before a payment is released, the financial consequence may be negligible.

If the same access persists for several hours, the potential consequences are much larger.

A trading-system compromise detected before the market opens is different from one discovered after a full day of activity.

A corrupted valuation file identified before NAV publication is different from one discovered after client statements are distributed.

A useful conceptual relationship is:

\begin{equation}
L_{\text{incident}}
=
f(
T_{\text{undetected}},
T_{\text{containment}},
A_{\text{authority}},
C_{\text{asset}},
M_{\text{market conditions}}
),
\end{equation}

where \(T_{\text{undetected}}\) is the period before detection, \(T_{\text{containment}}\) is the time required to restrict the event, \(A_{\text{authority}}\) represents the authority available to the compromised identity or agent, \(C_{\text{asset}}\) represents business criticality, and \(M_{\text{market conditions}}\) captures the financial environment.

The relationship is not necessarily linear. A short delay during a quiet period may have limited effect. The same delay during a market crisis may produce significant losses.

For this reason, mean time to detect and mean time to respond should not be viewed as abstract technology metrics. They are drivers of operational-loss severity.

\section{What the SOC Actually Observes}

The SOC normally observes events generated by many systems:

\begin{itemize}
    \item human and machine authentication;
    \item endpoint activity;
    \item cloud control planes;
    \item network connections;
    \item application logs;
    \item payment systems;
    \item trading infrastructure;
    \item databases;
    \item software-development systems;
    \item file-integrity monitors;
    \item third-party services;
    \item AI agents and autonomous workflows.
\end{itemize}

No individual log normally provides the full story.

The analyst must correlate information across sources.

For example, suppose the institution sees:

\begin{enumerate}
    \item a login from a new device;
    \item access to a sensitive database;
    \item a sudden export of records;
    \item a change to a privileged configuration.
\end{enumerate}

Each event has meaning.

Together they form a stronger pattern.

This can be expressed conceptually as:

\begin{equation}
\text{Signal Strength}
\neq
\sum \text{Individual Events}.
\end{equation}

The relationship among events matters.

\section{Rule-Based Detection}

Traditional detection often begins with explicit rules.

A rule may say:

\begin{quote}
Generate an alert when a privileged account produces more than five failed authentication events within ten minutes.
\end{quote}

Or:

\begin{quote}
Generate an alert when a production valuation file changes outside the approved deployment process.
\end{quote}

Rules have several advantages.

They are transparent.

They are testable.

They are easy to explain to management and auditors.

They are useful for clearly prohibited activity.

But rules are limited. A careful adversary may stay below a threshold. A legitimate operational process may trigger the rule. New attack patterns may not resemble historical signatures.

Therefore, mature monitoring combines deterministic rules with broader behavioral analysis.

\section{Behavioral and Anomaly Detection}

Behavioral detection asks whether current activity differs materially from expected activity.

Suppose a portfolio administrator normally logs in from Mexico City between 7:00 a.m. and 7:00 p.m., accesses three internal systems, and transfers only small files.

Now the same account begins accessing an unfamiliar system, at an unusual time, from a new device, while moving much more data than normal.

No single action needs to be prohibited.

The pattern creates concern.

A simplified anomaly score can be written as:

\begin{equation}
Z_t =
w_1 D_{\text{identity}}
+
w_2 D_{\text{device}}
+
w_3 D_{\text{time}}
+
w_4 D_{\text{resource}}
+
w_5 D_{\text{volume}},
\end{equation}

where each \(D\) measures deviation from an expected baseline.

The challenge is false positives.

Financial institutions are dynamic. Employees travel. Markets open overnight. Emergency activity occurs. New systems are introduced. A good detection program must adapt without becoming permissive.

\section{The AI Threat Lens: The Capability Frontier Is Moving}

The most important change in cybersecurity is not simply that AI can write code. It is that frontier models are becoming capable of sustained security analysis over longer and more complex tasks.

OpenAI states that GPT-5.6 is more capable in cybersecurity than its earlier models and that its testing currently shows greater strength in finding and fixing vulnerabilities than in carrying out reliable autonomous end-to-end attacks against hardened targets.\footnote{\url{https://openai.com/index/gpt-5-6/}}

OpenAI's deployment-safety documentation similarly reports that GPT-5.6 Sol and Terra can identify vulnerabilities and components of exploits, while not yet reliably completing autonomous attacks against hardened targets.\footnote{\url{https://deploymentsafety.openai.com/gpt-5-6}}

Anthropic's Mythos work shows an even more striking shift in vulnerability research. Anthropic reports that Mythos Preview identified previously unknown vulnerabilities across major operating systems and browsers and has used Mythos-class capabilities in Project Glasswing to find and help remediate weaknesses in critical software.\footnote{\url{https://www.anthropic.com/research/mythos-preview}}\footnote{\url{https://www.anthropic.com/glasswing}}

The UK AI Security Institute has also reported rapid improvements in autonomous cyber capability. Its 2026 work indicates that the reliable time horizon over which frontier models can perform cyber tasks has been increasing quickly, and that tool access and better scaffolding materially improve performance.\footnote{\url{https://www.aisi.gov.uk/blog/how-fast-is-autonomous-ai-cyber-capability-advancing}}\footnote{\url{https://www.aisi.gov.uk/frontier-ai-trends-report}}

For financial institutions, the implication is not that AI has made defense obsolete.

The implication is that historical assumptions about the scarcity, speed, and persistence of expert cyber labor are changing.

\section{AI and Vulnerability Discovery}

Vulnerability discovery is one of the areas where frontier AI is advancing fastest.

A vulnerability is a weakness in software or system design that can create an unintended security consequence.

Historically, finding subtle vulnerabilities often required highly skilled humans who could understand large code bases, reason across complex program states, and test unusual interactions.

AI changes the economics of that work.

A capable model can assist with:

\begin{itemize}
    \item understanding unfamiliar source code;
    \item tracing data and control flows;
    \item comparing implementation with intended behavior;
    \item identifying suspicious code patterns;
    \item reviewing patches;
    \item summarizing possible impact;
    \item prioritizing areas for human investigation.
\end{itemize}

OpenAI's trusted-access cybersecurity work describes frontier models assisting with understanding unfamiliar code, mapping affected surfaces, tracing root causes, reviewing patches, and producing remediation guidance.\footnote{\url{https://openai.com/index/gpt-5-5-with-trusted-access-for-cyber/}}

Anthropic's coordinated vulnerability disclosure effort similarly uses human review to validate AI-discovered findings before disclosure to maintainers.\footnote{\url{https://red.anthropic.com/2026/cvd/}}

The defensive opportunity is substantial.

The threat is that similar analytical efficiency can reduce the time required for malicious actors to identify weak systems.

\begin{risklens}
The financial consequence of faster vulnerability discovery is a shorter defensive window. Institutions may have less time between public disclosure, automated analysis, and attempted misuse. Patch governance therefore becomes economically more important.
\end{risklens}

\section{Typical AI-Enabled Threat Tactics}

The following tactics are described only at the level required for threat modeling and governance.

\subsection{1. Automated Surface Analysis}

AI systems can process large inventories of software, configurations, documents, and system descriptions more quickly than human teams.

The strategic effect is breadth.

A threat actor can potentially investigate more systems with the same amount of human supervision.

The defensive response is accurate asset inventory, exposure management, and prioritization of critical systems.

\subsection{2. Automated Vulnerability Triage}

AI can rank suspected weaknesses, compare them with system context, and focus attention on the most promising findings.

The strategic effect is efficiency.

The defensive response is to use similar capability internally so that defenders identify and patch weaknesses first.

\subsection{3. Long-Horizon Reasoning}

Frontier models increasingly sustain work over longer sequences.

Rather than producing one isolated answer, an agent can maintain an objective, evaluate results, revise its plan, and continue.

The strategic effect is persistence.

The defensive response is continuous monitoring rather than reliance on isolated point-in-time controls.

\subsection{4. Parallel Agent Teams}

A supervisor may assign different analytical tasks to specialized agents.

At a high level:

\begin{equation}
\text{Supervisor}
\rightarrow
\begin{cases}
\text{Code Analysis Agent}\\
\text{Identity Analysis Agent}\\
\text{Configuration Agent}\\
\text{Evidence Agent}
\end{cases}
\rightarrow
\text{Coordinator}.
\end{equation}

The strategic effect is parallelism.

The defensive response is also parallelism: automated defensive review, continuous scanning, and AI-supported SOC triage.

\subsection{5. Adaptive Workflow Selection}

A conventional script follows predetermined steps.

An AI agent can choose among alternative actions based on observed results.

The strategic effect is adaptability.

This makes behavior less predictable and increases the value of detecting outcomes and anomalies rather than only known signatures.

\subsection{6. Memory-Augmented Persistence}

An agent can maintain state across tasks.

It may remember prior findings, rejected hypotheses, system relationships, and observed behavior.

The strategic effect is cumulative learning.

Defenders should assume that repeated low-level probing may be connected by persistent machine memory even when individual events appear weak.

\subsection{7. AI-Assisted Social Engineering}

AI can produce highly fluent, context-aware communications.

The relevant risk for finance is not literary quality. It is the possibility that fraudulent requests become more personalized and operationally convincing.

This increases the importance of independent verification, dual authorization, and transaction controls.

\subsection{8. Tool Chaining}

An agent can combine several legitimate capabilities into one workflow.

For example, one tool may retrieve information, another may analyze it, and another may create a workflow action.

The strategic effect is composability.

The defensive response is to constrain tool permissions and prevent untrusted information from automatically producing high-impact actions.

\section{AI Does Not Need to Be Perfect to Increase Threat}

This point is important.

A model does not need to autonomously compromise the most hardened bank in the world to materially change cyber risk.

It only needs to reduce the cost of some portion of the attack process.

If AI makes reconnaissance faster, vulnerability analysis cheaper, social engineering more scalable, or repeated experimentation more persistent, the economics have changed.

This can be represented as:

\begin{equation}
C_{\text{attack}}
=
C_{\text{expertise}}
+
C_{\text{time}}
+
C_{\text{coordination}}
+
C_{\text{computation}}.
\end{equation}

AI can reduce the first three components.

As the cost falls, more actors may be able to perform work that previously required specialized teams.

\section{The Defensive Opportunity: AI in the SOC}

The same capabilities can benefit defenders.

A defensive AI system can summarize a large alert queue, reconstruct timelines, correlate identity and endpoint evidence, explain why an alert matters, identify similar cases, and recommend investigation steps.

A useful architecture is:

\begin{equation}
\text{Telemetry}
\rightarrow
\text{Deterministic Detection}
\rightarrow
\text{AI Correlation}
\rightarrow
\text{Risk Prioritization}
\rightarrow
\text{Human Decision}.
\end{equation}

AI should initially be strongest in interpretation rather than irreversible action.

For example, an AI agent may recommend disabling an account. Whether it can actually disable the account should depend on the institution's risk architecture.

\section{Typical Defensive Agent Architecture}

A mature financial SOC may eventually operate several specialized defensive agents.

\subsection{Identity Agent}

The identity agent monitors authentication, privileges, devices, and unusual account behavior.

\subsection{Endpoint Agent}

The endpoint agent reviews activity on workstations and servers.

\subsection{Data Integrity Agent}

The integrity agent monitors changes to critical files, data sets, and configurations.

\subsection{Vulnerability Agent}

The vulnerability agent reviews software inventories, patch status, known weaknesses, and AI-assisted code findings.

\subsection{Financial Context Agent}

This agent maps technical assets to financial processes.

It can answer questions such as:

\begin{quote}
Does this server support payment release, NAV calculation, trade execution, or regulatory reporting?
\end{quote}

\subsection{Supervisor Agent}

The supervisor combines findings and creates a consolidated case.

This architecture can be expressed as:

\begin{equation}
\left\{
A_I,
A_E,
A_D,
A_V,
A_F
\right\}
\rightarrow
A_S
\rightarrow
\text{Human Incident Lead}.
\end{equation}

The architecture is powerful because no single agent needs to understand the entire institution.

But it creates governance risk.

If the supervisor receives incorrect information from one specialist agent, the error can propagate.

Agent outputs therefore require provenance, confidence, and independent evidence.

\section{Deterministic Controls Must Remain}

AI is useful when the problem requires interpretation.

It is less appropriate when the institution already knows the rule.

A payment above an approved limit should require the appropriate approval.

A trader should not exceed a position limit because an AI agent believes market conditions justify it.

A production server should not accept an unapproved software deployment.

A restricted security should not be traded by an unauthorized portfolio.

These rules should remain deterministic.

\begin{intuition}
Use AI for ambiguity. Use deterministic controls for policy.
\end{intuition}

\section{Detection Engineering in the AI Era}

Detection engineering is the discipline of designing, testing, and improving rules that identify suspicious behavior.

AI changes detection engineering in several ways.

First, AI can help analysts explain large rule sets.

Second, it can compare alerts with historical incidents.

Third, it can propose candidate correlations.

Fourth, it can search large documentation sets for relevant context.

Fifth, it can assist defenders in reviewing software and configurations before an attacker finds the weakness.

But AI-generated detections should not automatically become production controls.

They require validation.

A poorly designed detection can create operational disruption.

The appropriate process is:

\begin{equation}
\text{AI Proposal}
\rightarrow
\text{Human Review}
\rightarrow
\text{Historical Testing}
\rightarrow
\text{Controlled Deployment}
\rightarrow
\text{Performance Monitoring}.
\end{equation}

This resembles model validation in finance.

\section{Patch Management Becomes a Strategic Control}

As AI accelerates vulnerability discovery, patch management becomes more important.

Traditional patching sometimes treated every update as a routine IT task.

For a financial institution, prioritization should depend on:

\begin{itemize}
    \item exposure;
    \item exploitability;
    \item business criticality;
    \item privilege;
    \item data sensitivity;
    \item availability of compensating controls;
    \item operational cost of patching.
\end{itemize}

A vulnerability in an isolated test system is different from one in an internet-facing payment gateway.

The institution needs a risk-based patching framework.

A simplified priority score can be written as:

\begin{equation}
P_{\text{patch}}
=
V
\times
E
\times
C
\times
A,
\end{equation}

where \(V\) is technical severity, \(E\) exposure, \(C\) business criticality, and \(A\) current adversarial activity or likelihood.

The formula is conceptual.

Its purpose is to prevent technical severity from being interpreted without financial context.

\section{Protection Strategy: Shrink the Attack Surface}

The first defense against AI-accelerated vulnerability discovery is to reduce the number of unnecessary things that can be found.

Financial institutions should:

\begin{itemize}
    \item maintain accurate asset inventories;
    \item retire unsupported software;
    \item remove unnecessary public services;
    \item disable unused accounts;
    \item minimize exposed administrative interfaces;
    \item segment critical systems;
    \item restrict machine and agent permissions;
    \item document ownership.
\end{itemize}

AI may make discovery faster.

It cannot discover an unnecessary service that has been removed.

\section{Protection Strategy: Defender-First Vulnerability Discovery}

Institutions should increasingly use advanced AI defensively.

The objective is simple:

\begin{equation}
T_{\text{defender discovery}}
<
T_{\text{attacker discovery}}.
\end{equation}

Internal code review, AI-assisted vulnerability research, automated configuration review, and continuous testing can help identify weaknesses before external actors do.

The institution must govern this work carefully, particularly when advanced cyber models or tools are involved.

The goal is remediation, not uncontrolled experimentation.

\section{Protection Strategy: Monitor the Agent Layer}

As organizations deploy internal agents, the SOC must monitor them as a new class of identity.

Each agent should produce telemetry on:

\begin{itemize}
    \item model identity and version;
    \item human or business owner;
    \item tools used;
    \item data accessed;
    \item actions proposed;
    \item actions executed;
    \item approval events;
    \item failures;
    \item policy exceptions.
\end{itemize}

An autonomous agent without observability is an operational-risk blind spot.

\section{Protection Strategy: Limit Autonomous Response}

Security teams are attracted to automated containment because speed matters.

But automatic response can itself create financial loss.

Suppose an AI defender incorrectly concludes that a legitimate trading gateway is compromised and disables it during a volatile market.

The defensive system has created an operational incident.

Therefore, autonomous response should be tiered.

Low-impact reversible actions may be automated.

High-impact actions should require stronger evidence or human approval.

A useful hierarchy is:

\begin{equation}
\text{Observe}
<
\text{Recommend}
<
\text{Restrict}
<
\text{Disable}
<
\text{Destroy or Rebuild}.
\end{equation}

Authority should decrease as reversibility decreases.

\section{Financial Practice: A Payment SOC}

Consider a fictional bank whose SOC receives an alert involving the payment platform.

The evidence shows:

\begin{enumerate}
    \item a privileged account authenticated from a new device;
    \item the account viewed several beneficiary records;
    \item no payment was released;
    \item the activity occurred shortly before the daily high-value payment window.
\end{enumerate}

A conventional SOC may classify the alert based on identity behavior.

An AI-assisted SOC can enrich the case with business context:

\begin{itemize}
    \item the account can prepare but not release payments;
    \item the employee is traveling;
    \item multifactor authentication succeeded;
    \item the beneficiary records were part of a scheduled review;
    \item no approval limits changed.
\end{itemize}

The enriched evidence may reduce the severity.

This is a good use of AI: interpreting context.

Now suppose an autonomous defensive agent wants to disable the account immediately.

Management must consider the cost of false action. If the account is needed for a critical payment window, unnecessary disabling may create a liquidity or client-service problem.

The best architecture combines rapid AI analysis with proportionate controls.

\section{Financial Practice: An Exchange}

An exchange has a different risk profile.

Availability is central.

A false positive that disconnects a participant may disrupt trading.

A false negative that misses an attack on the matching or market-data environment may have broader market consequences.

The SOC must therefore combine technical severity with market criticality.

During volatile markets, escalation thresholds may need to change.

Incident response should be coordinated with market operations, surveillance, clearing, communications, and regulators.

This demonstrates why cybersecurity cannot be isolated from the financial operating model.

\section{Mini Case: AI-Discovered Vulnerability in a Pricing Service}

A fictional asset manager operates an internally developed pricing service.

A defensive AI review identifies a plausible software weakness in a component used to import pricing data.

The model assigns high confidence.

What should the firm do?

It should not automatically assume the finding is correct.

The appropriate process is:

\begin{enumerate}
    \item human validation of the finding;
    \item identification of affected systems;
    \item assessment of business criticality;
    \item evaluation of whether the service is externally exposed;
    \item development and review of a patch;
    \item testing in a controlled environment;
    \item deployment through approved change management;
    \item post-deployment verification;
    \item review of historical logs for evidence of prior misuse.
\end{enumerate}

The important point is that AI accelerates discovery.

Governance still determines remediation.

\begin{risklens}
AI may shorten the time required to find a vulnerability, but a financial institution should not shorten its control discipline. Speed in discovery must be matched by speed in governed validation and remediation.
\end{risklens}

\section{Safe Notebook Lab}

\begin{notebooklab}
The companion notebook should create a synthetic SOC for a fictional financial institution.

Students can generate harmless events representing:
\begin{enumerate}
    \item normal authentication;
    \item unusual authentication;
    \item a synthetic file-integrity change;
    \item a fictional vulnerability finding;
    \item an agent action;
    \item a payment-related alert.
\end{enumerate}

The notebook can include several transparent defensive agents:
\begin{itemize}
    \item Identity Agent;
    \item Integrity Agent;
    \item Vulnerability Agent;
    \item Financial Context Agent;
    \item Supervisor Agent.
\end{itemize}

Each agent should operate only on local synthetic records.

Students should compare:
\begin{enumerate}
    \item deterministic rules;
    \item AI-style risk scoring;
    \item human escalation decisions.
\end{enumerate}

A second exercise should model patch prioritization using synthetic severity, exposure, criticality, and compensating-control variables.

The notebook must not scan external systems, test public services, generate operational exploitation instructions, use real credentials, or connect autonomous agents to real infrastructure.
\end{notebooklab}

\section{A Pragmatic SOC Checklist for Finance Professionals}

Senior management and risk committees should ask:

\begin{enumerate}
    \item Which financial processes are mapped to critical technical assets?
    \item Are identity, endpoint, cloud, application, and agent events centrally observable?
    \item How quickly are high-impact alerts reviewed?
    \item Which alerts can produce automatic actions?
    \item What is the business cost of a false positive?
    \item What is the expected loss from a false negative?
    \item How are AI-discovered vulnerabilities validated?
    \item How quickly are critical vulnerabilities remediated?
    \item Does the institution use advanced AI defensively before adversaries do?
    \item Are internal AI agents included in SOC monitoring?
    \item Can one defensive agent disable a critical financial service without independent approval?
    \item Can management reconstruct exactly what happened after an incident?
\end{enumerate}

\section{Chapter Summary}

The main lessons of this chapter are:

\begin{enumerate}
    \item Security operations are a financial control function because detection speed and containment directly affect loss severity.
    \item Effective monitoring requires collection, normalization, business enrichment, detection, triage, escalation, and response.
    \item Rule-based detection remains essential, but behavioral analysis is needed for patterns that do not violate one explicit rule.
    \item Frontier AI models are rapidly improving at vulnerability discovery, code analysis, and sustained cyber tasks.
    \item AI changes cyber risk even before fully autonomous end-to-end attacks become consistently reliable, because it reduces the cost of reconnaissance, analysis, experimentation, coordination, and persistence.
    \item Typical AI-enabled tactics include automated surface analysis, vulnerability triage, long-horizon reasoning, parallel agent teams, adaptive workflows, persistent memory, social engineering, and tool chaining.
    \item The same AI capabilities provide a major defensive opportunity through automated code review, alert correlation, vulnerability discovery, and SOC triage.
    \item Patch management becomes more strategically important as AI reduces the time needed to identify weaknesses.
    \item Defensive agents must be monitored and constrained because an incorrect automated response can itself create financial loss.
    \item The objective is not fully autonomous security. It is governed, observable, financially informed defensive automation.
\end{enumerate}

\section{Review Questions}

\begin{enumerate}
    \item Why are mean time to detect and mean time to respond financial-risk variables?
    \item What is the difference between rule-based detection and behavioral detection?
    \item Why does business criticality change the meaning of a technical alert?
    \item How does frontier AI change the economics of vulnerability discovery?
    \item Why does AI not need to execute a complete autonomous attack to increase cyber risk?
    \item Describe three high-level AI-enabled tactics that increase scale or persistence.
    \item Why can multi-agent architectures create a greater monitoring challenge than a single model?
    \item How can AI strengthen a security operations center?
    \item Why should AI-generated vulnerability findings be validated before remediation?
    \item What does defender-first vulnerability discovery mean?
    \item Why can autonomous defensive response create operational risk?
    \item Which actions should remain deterministic or human-approved in a financial SOC?
\end{enumerate}
